\documentclass[11pt,a4paper]{article}

\usepackage[utf8]{inputenc}
\usepackage[T1]{fontenc}
\usepackage{amsmath,amssymb,amsfonts}
\usepackage{bm}
\usepackage{physics}
\usepackage{geometry}
\usepackage{enumitem}
\usepackage{booktabs}
\usepackage{hyperref}
\usepackage{cleveref}

\geometry{margin=2.5cm}

\newcommand{\avg}[1]{\langle #1 \rangle}
\newcommand{\avgJ}[1]{\overline{#1}}
\newcommand{\re}{\operatorname{Re}}

\title{%
  Validit\`a dell'approssimazione gaussiana i.i.d. \\[4pt]
  per gli accoppiamenti $g^{(2)}$ e $g^{(4)}$ \\[4pt]
  \large Modello $p$-spin $(2{+}4)$ per random laser:
         caso denso e sparso}
\author{}
\date{\today}

\begin{document}
\maketitle

%======================================================================
\section{Origine fisica degli accoppiamenti}\label{sec:origin}
%======================================================================

\subsection{Accoppiamenti a 4 corpi: $\chi^{(3)}$}

Nell'espansione dei modi normali del campo elettromagnetico in un mezzo
disordinato con suscettivit\`a non lineare del terzo ordine,
l'accoppiamento quartico tra i modi $(i_1, i_2, i_3, i_4)$ \`e
(Eq.~(8) di \cite{Gradenigo2020}, e \cite{Antenucci2015a}):
\begin{equation}\label{eq:g4_integral}
  g^{(4)}_{i_1 i_2 i_3 i_4}
  \propto
  \int_V d\vb{r}\;
  \sum_{\nu}
    E^{(\nu_1)}_{i_1}(\vb{r})\,
    E^{(\nu_2)}_{i_2}(\vb{r})\,
    E^{(\nu_3)}_{i_3}(\vb{r})\,
    E^{(\nu_4)}_{i_4}(\vb{r})\;
    \chi^{(3)}_{\nu}\!\bigl(\vb{r}\,\big|\,
      \omega_{i_1},\omega_{i_2},\omega_{i_3},\omega_{i_4}\bigr).
\end{equation}
L'integrale \`e esteso al volume del mezzo $V$, e
$\chi^{(3)}_\nu(\vb{r})$ \`e la suscettivit\`a non lineare del terzo
ordine, spazialmente disordinata e dipendente dalle frequenze.

\subsection{Accoppiamenti a 2 corpi: $\chi^{(1)}$}

Il termine bilineare nell'Hamiltoniana proviene dalla polarizzazione
lineare del mezzo, mediata dalla suscettivit\`a del primo ordine
$\chi^{(1)}$. L'accoppiamento fra i modi $i$ e $j$ ha la forma
(\cite{Angelani2006, Leuzzi2009, Antenucci2015a}):
\begin{equation}\label{eq:g2_integral}
  g^{(2)}_{ij}
  \propto
  \int_V d\vb{r}\;
  \sum_\nu
    E^{(\nu)}_i(\vb{r})\,
    E^{(\nu)*}_j(\vb{r})\;
    \chi^{(1)}_\nu\!\bigl(\vb{r}\,\big|\,\omega_i,\omega_j\bigr).
\end{equation}
Questa espressione ha una struttura analoga all'\cref{eq:g4_integral},
con due profili di modo invece di quattro.

\paragraph{Struttura del tensore $g^{(2)}$.}
La parte diagonale $g^{(2)}_{ii} = g(\omega_i)$ coincide con il
\emph{profilo di guadagno netto} (net gain): guadagno del mezzo attivo
meno le perdite della cavit\`a aperta. Nell'articolo di Gradenigo
et~al.\@ \cite{Gradenigo2020}, questa parte \`e posta $g_{ii} = g =
\mathrm{const}$. Con il vincolo sferico
$\sum_k |a_k|^2 = \epsilon N$ il termine diagonale diventa una costante
$\mathcal{H}_2^{\mathrm{diag}} = -g\,\epsilon N$ e si disaccoppia dalla
dinamica.

La parte \emph{off-diagonale} ($i \neq j$) descrive l'accoppiamento
lineare tra modi della cavit\`a passiva, mediato dalla loro
sovrapposizione spaziale nel mezzo disordinato
\cite{Hackenbroich2003, Viviescas2003}. Questo termine \`e significativo
solo quando le frequenze dei due modi sono quasi-degeneri
$|\omega_i - \omega_j| \lesssim \gamma$ (condizione di quasi-risonanza,
FMC per il termine a 2~corpi).

%======================================================================
\section{Argomento per la gaussianit\`a}\label{sec:clt}
%======================================================================

Tanto l'\cref{eq:g4_integral} quanto l'\cref{eq:g2_integral} sono
integrali spaziali su quantit\`a microscopiche --- i profili di modo
$E_i(\vb{r})$ e la suscettivit\`a $\chi^{(p)}(\vb{r})$ --- che in
un mezzo random sono esse stesse variabili aleatorie con forti
fluttuazioni spaziali \cite{Fallert2009, Conti2008}.

Si pu\`o decomporre il volume in celle $\{\Delta V_\ell\}$ di dimensione
paragonabile alla lunghezza di correlazione spaziale delle fluttuazioni:
\begin{equation}\label{eq:g4_sum}
  g^{(4)}_{i_1 i_2 i_3 i_4}
  \approx
  \sum_\ell \xi_\ell^{(4)}(i_1,i_2,i_3,i_4)\,,
\end{equation}
dove $\xi_\ell^{(4)}$ \`e il contributo della cella $\ell$-esima.
Analogamente per $g^{(2)}$:
\begin{equation}\label{eq:g2_sum}
  g^{(2)}_{ij}
  \approx
  \sum_\ell \xi_\ell^{(2)}(i,j)\,.
\end{equation}

Quando il volume del mezzo contiene molte ($\gg 1$) lunghezze di
correlazione, i contributi $\xi_\ell$ sono \emph{debolmente correlati}
(solo celle adiacenti hanno correlazione sensibile), e il teorema del
limite centrale (CLT) implica che ciascun coupling converge in
distribuzione a una gaussiana:
\begin{equation}\label{eq:gaussianity}
  g^{(p)}
  \;\xrightarrow{\;V \gg \xi^d\;}
  \;\mathcal{N}\!\left(\mu^{(p)},\; [\sigma^{(p)}]^2\right).
\end{equation}
L'approssimazione \`e tanto migliore quanto pi\`u grande \`e il rapporto
$V / \xi^d$ (volume del campione su volume di correlazione spaziale).
Gradenigo~et~al.\@ la adottano esplicitamente in Eq.~(A4) di
\cite{Gradenigo2020}.

\paragraph{Nota.}
L'argomento CLT giustifica la gaussianit\`a della distribuzione
\emph{marginale} di ciascun coupling, ma \textbf{non dice nulla}
sull'indipendenza tra couplings diversi. La questione dell'indipendenza
richiede un'analisi separata (\cref{sec:iid}).

%======================================================================
\section{Indipendenza (o meno) degli accoppiamenti}\label{sec:iid}
%======================================================================

\subsection{Struttura delle correlazioni}

Due accoppiamenti $g^{(4)}_{\bar\imath}$ e $g^{(4)}_{\bar\jmath}$ che
condividono $k$ indici di modo ($0 \le k \le 3$) sono in generale
\emph{correlati}: gli integrali di sovrapposizione coinvolgono gli
stessi profili di modo globali $E_i(\vb{r})$. Il momento misto
$\avgJ{g^{(4)}_{\bar\imath}\, g^{(4)}_{\bar\jmath}} -
 \avgJ{g^{(4)}_{\bar\imath}}\,\avgJ{g^{(4)}_{\bar\jmath}}$
\`e in generale non nullo quando $\bar\imath$ e $\bar\jmath$ hanno
almeno un modo in comune.

Per i $g^{(2)}$, analogamente, $g^{(2)}_{ij}$ e $g^{(2)}_{ik}$
($j \neq k$) sono correlati perch\'e condividono il profilo
$E_i(\vb{r})$.

Questa struttura di correlazione \`e discussa da Zaitsev e
Deych~\cite{Zaitsev2010}.

\subsection{Caso denso: perch\'e le correlazioni sono trascurabili}

\paragraph{Grafo completo ($\binom{N}{4}$ quartetti).}
Ogni modo $i$ partecipa a $\binom{N-1}{3} = O(N^3)$ quartetti.
La correlazione tra $g^{(4)}_{\bar\imath}$ e $g^{(4)}_{\bar\jmath}$
che condividono il modo $i$ decade come $1/V$ (dove $V$ \`e il volume
del mezzo), perch\'e i tre profili di modo \emph{non condivisi} hanno
overlap spaziale trascurabile.

In termini di contributo alla funzione di partizione: nella somma
$H_4 = \sum_{\bar\imath} g^{(4)}_{\bar\imath} (\cdots)$, gli
$O(N^4)$ termini autocorrelati formano una frazione $O(1/N)$ del totale
$O(N^4) \times O(N^4)$ delle coppie di couplings. Per la legge dei
grandi numeri le correlazioni sono media\-mente trascurabili nella
termodinamica: le osservabili termodinamiche dipendono solo dai primi
due momenti della distribuzione congiunta, che per un sistema con
$O(N^4)$ termini indipendenti sono ben approssimati dai momenti
marginali.

L'argomento \`e analogo per $g^{(2)}$: ci sono $N(N-1)/2$ coppie, e
le correlazioni fra coppie che condividono un indice pesano $O(1/N)$
nella somma.

\paragraph{Grafo FMC (mode-locked, $O(N^3)$ quartetti).}
La frequency-matching condition riduce il numero di quartetti da
$O(N^4)$ a $O(N^3)$ \cite{Marruzzo2018}. Ogni modo partecipa ancora a
$O(N^2)$ quartetti. L'argomento precedente si applica con la stessa
struttura: la correlazione media fra coppie di couplings che condividono
un modo scompare nel limite termodinamico come $O(1/N)$.

Per i $g^{(2)}$, la FMC seleziona solo le coppie con
$|\omega_i - \omega_j| \le \gamma$. Con frequenze uniformi su un
intervallo, il numero di coppie sopravvissute \`e $O(N)$ (ciascun modo
si accoppia a $O(1)$ vicini in frequenza). In questo regime il termine
$H_2$ ha una connettivit\`a \emph{finita}, e $g^{(2)}_{ij}$ e
$g^{(2)}_{ik}$ sono correlati con peso $O(1)$.
Tuttavia, poich\'e il termine $H_2$ \`e bilineare, la correlazione
tra coppie non modifica la classe di universalit\`a della transizione
vetrosa (il termine quadratico aggiunge un campo esterno effettivo,
non un'interazione aggiuntiva nel senso di replica).

\bigskip
\noindent
\textbf{Conclusione per il caso denso.}
L'approssimazione gaussiana i.i.d.\@ \`e pienamente giustificata sia
per $g^{(2)}$ che per $g^{(4)}$. La gaussianit\`a segue dal CLT
(\cref{sec:clt}); l'indipendenza segue dalla soppressione delle
correlazioni nel limite $N \to \infty$.

%======================================================================
\section{Caso sparso: $O(N)$ quartetti
         sub-campionati}\label{sec:sparse}
%======================================================================

Nella variante sparsa implementata nel codice, si selezionano
(via Fisher--Yates parziale) $n_s = N$ quartetti uniformemente a caso
dall'insieme dei quartetti FMC-sopravvissuti ($O(N^3)$), ottenendo
\begin{equation}\label{eq:H4_sparse}
  H_4^{\mathrm{sparse}}
  = -\sum_{\mu=1}^{N} \re\!\bigl[
      \tilde g^{(4)}_\mu\, f_{c_\mu}(a_{i_\mu},a_{j_\mu},a_{k_\mu},a_{l_\mu})
    \bigr],
\end{equation}
con nuova varianza $\mathrm{Var}(\tilde g^{(4)}) = J^2 N / n_s = J^2$
per garantire estensivit\`a (Eq.~(A4) generalizzata).

\subsection{Gaussianit\`a: ancora valida}

L'argomento CLT della \cref{sec:clt} \`e indipendente dalla topologia
della rete di interazione: riguarda la distribuzione del singolo
coupling, che \`e un integrale spaziale su molti contributi debolmente
correlati. Il sub-campionamento seleziona quartetti in modo
\emph{indipendente dal valore} del coupling (Fisher--Yates su indici,
non su valori), quindi non introduce bias sulla distribuzione marginale.
\begin{equation}
  \tilde g^{(4)}_\mu
  \sim \mathcal{N}(0,\, J^2) \qquad \forall\, \mu = 1,\ldots,N.
\end{equation}

Lo stesso argomento si applica ai $g^{(2)}$ sopravvissuti dopo il
filtro FMC: il filtro agisce sulle frequenze, che sono indipendenti
dalla realizzazione dei couplings.

\subsection{Indipendenza: pi\`u delicata}\label{sec:sparse_iid}

Con $N$ quartetti e $N$ modi, ogni modo partecipa in media a
$4N / N = 4$ quartetti. La connettivit\`a di ciascun modo \`e dunque
$O(1)$, come in un vetro di spin su grafo sparso tipo
Viana--Bray~\cite{Viana1985}.

In questo regime:
\begin{enumerate}[nosep]
  \item \textbf{Le correlazioni non sono soppresse da $N$.}
    Due couplings che condividono un modo hanno una correlazione $O(1/V)$
    (dal contributo spaziale), ma la loro vicinanza nella rete \emph{non}
    \`e diluita da un fattore $1/N$: un dato modo vede solo $O(1)$
    quartetti, non $O(N^2)$. L'effetto relativo delle correlazioni sulla
    termodinamica \emph{locale} (campo molecolare su un singolo modo)
    \`e quindi pi\`u pronunciato.

  \item \textbf{L'universalit\`a rispetto alla forma della distribuzione
    non \`e garantita a priori.}
    Su grafi densi, il teorema della cavit\`a (o il metodo delle
    repliche con $O(N^p)$ termini) mostra che solo i primi due momenti
    della distribuzione dei couplings entrano nel calcolo al punto sella.
    Su grafi sparsi con connettivit\`a finita, i momenti superiori della
    distribuzione possono influire sulle propriet\`a termodinamiche:
    ad esempio, la temperatura critica di un modello $\pm J$ su grafo
    di Bethe non coincide in generale con quella di un modello
    gaussiano con la stessa varianza~\cite{MPV1987}.

  \item \textbf{L'auto-mediazione (self-averaging) \`e pi\`u debole.}
    Con $O(N)$ termini nella somma dell'Hamiltoniana, la fluttuazione
    relativa dell'energia per campione di disordine decade come
    $1/\sqrt{N}$, pi\`u lentamente che nel caso denso ($O(N^3)$ o
    $O(N^4)$ termini).
\end{enumerate}

\subsection{Argomenti a favore dell'approssimazione gaussiana i.i.d.\@
            nel caso sparso}

Nonostante le cautele della sezione precedente, l'approssimazione
i.i.d.\@ gaussiana resta ragionevole nelle simulazioni per le seguenti
ragioni:

\begin{enumerate}[nosep]
  \item \textbf{Il sub-campionamento \`e uniforme e indipendente.}
    I quartetti sono selezionati uniformemente da un pool di $O(N^3)$
    candidati. La probabilit\`a che due quartetti selezionati
    condividano un modo \`e:
    \begin{equation}
      P(\text{modo condiviso})
      = 1 - \prod_{k=1}^{3}\frac{\binom{N-4}{4-k}\binom{N}{k}}{\binom{N}{4}-k+1}
      \approx \frac{16}{N}\,,
    \end{equation}
    che per $N \gg 1$ \`e trascurabile. La maggior parte delle coppie di
    couplings selezionati coinvolge indici completamente disgiunti, e
    dunque \`e automaticamente indipendente (anche senza invocare la
    decorrelazione spaziale).

  \item \textbf{Correlazione spaziale $O(1/V)$.}
    Anche quando due couplings condividono un indice, la correlazione
    proviene dall'overlap in $\vb{r}$ dei profili di modo non condivisi
    e decade come $1/V$. In un mezzo macroscopico questa contribuzione
    \`e trascurabile per qualunque topologia della rete.

  \item \textbf{Stabilit\`a della transizione.}
    La transizione vetrosa nel modello sferico $p$-spin \`e di tipo
    RFOT (random first-order transition) \cite{Kirkpatrick1989}.
    Lo scenario RFOT \`e stabile rispetto a perturbazioni della
    distribuzione dei couplings (es.\@ il modello ortogonale random
    \cite{Marinari1994, Parisi1995} ha diagramma di fase qualitativamente
    analogo pur avendo couplings fortemente correlati). Questo suggerisce
    che, per il comportamento critico, la forma gaussiana cattura la
    fisica essenziale anche con connettivit\`a finita.

  \item \textbf{Coerenza con la letteratura.}
    L'approssimazione gaussiana i.i.d.\@ \`e adottata esplicitamente in
    \cite{Gradenigo2020} (Eq.~A4), \cite{Antenucci2015a},
    \cite{Antenucci2015b}, e \cite{Angelani2006} per il caso
    FMC-denso. La sua estensione al caso sparso \`e un'assunzione di
    lavoro standard nel passaggio a grafi diluti di modelli
    mean-field~\cite{MPV1987}.
\end{enumerate}

%======================================================================
\section{Riepilogo}\label{sec:summary}
%======================================================================

\begin{center}
\renewcommand{\arraystretch}{1.3}
\begin{tabular}{lp{5.5cm}p{5.5cm}}
\toprule
 & \textbf{Gaussianit\`a} (CLT) & \textbf{Indipendenza} (i.i.d.) \\
\midrule
$g^{(2)}$ denso (FC)
  & \checkmark\; Integrale spaziale su $\gg 1$ celle di correlazione.
  & \checkmark\; Correlazioni fra coppie pesano $O(1/N)$ nella termodinamica. \\
$g^{(2)}$ FMC
  & \checkmark\; Stesso argomento CLT.
  & $\sim$\; Con $O(N)$ coppie sopravvissute, connettivit\`a finita. Termine bilineare; le correlazioni non alterano la classe di universalit\`a RFOT. \\
$g^{(4)}$ denso (FC o FMC)
  & \checkmark\; Eq.~(A4) di \cite{Gradenigo2020}.
  & \checkmark\; Ogni modo in $O(N^2)$ quartetti: decorrelazione $O(1/N)$. \\
$g^{(4)}$ sparso ($N$ quartetti)
  & \checkmark\; Il sub-campionamento non altera la distribuzione marginale.
  & $\triangle$\; Connettivit\`a $O(1)$; correlazioni non soppresse da $1/N$. Tuttavia: (a)~coppie selezionate condividono un modo con probabilit\`a $O(1/N)$; (b)~correlazione spaziale $O(1/V)$; (c)~lo scenario RFOT \`e stabile sotto perturbazioni della distribuzione. \\
\bottomrule
\end{tabular}
\end{center}

\bigskip
\noindent
\textbf{Conclusione.}
\begin{itemize}[nosep]
  \item La \textbf{gaussianit\`a} di ciascun coupling \`e ben giustificata
    in tutti i casi (denso e sparso, $g^{(2)}$ e $g^{(4)}$),
    a patto che il volume del mezzo contenga molte lunghezze di
    correlazione spaziale.
  \item L'\textbf{indipendenza} \`e rigorosamente giustificata solo nel
    limite termodinamico per reti dense ($O(N^3)$ o pi\`u termini per modo).
    Nel caso sparso con $O(N)$ quartetti e connettivit\`a $O(1)$,
    l'approssimazione i.i.d.\@ \`e un'ipotesi di lavoro \emph{non
    esattamente verificata}, ma supportata da:
    \begin{itemize}[nosep]
      \item la probabilit\`a $O(1/N)$ che due couplings sub-campionati
            condividano un modo;
      \item la decorrelazione spaziale $O(1/V)$;
      \item la robustezza dello scenario RFOT rispetto alla distribuzione
            dei couplings.
    \end{itemize}
  \item Per i $g^{(2)}$ con FMC, la connettivit\`a \`e gi\`a $O(1)$ anche
    senza sub-campionamento. L'approssimazione i.i.d.\@ \`e dunque
    paragonabile a quella del caso sparso $g^{(4)}$.
  \item Eventuali deviazioni sistematiche dal comportamento i.i.d.\@ sono
    attese essere pi\`u visibili a \emph{piccoli $N$} e nei
    \emph{momenti superiori} ($P(q)$, bimodalit\`a), dove la taglia
    finita del campione amplifica le fluttuazioni residue.
\end{itemize}

%======================================================================
\begin{thebibliography}{99}

\bibitem{Gradenigo2020}
  G.~Gradenigo, F.~Antenucci, and L.~Leuzzi,
  Phys.\ Rev.\ Research \textbf{2}, 023399 (2020).

\bibitem{Antenucci2015a}
  F.~Antenucci, C.~Conti, A.~Crisanti, and L.~Leuzzi,
  Phys.\ Rev.\ Lett.\ \textbf{114}, 043901 (2015).

\bibitem{Antenucci2015b}
  F.~Antenucci, A.~Crisanti, and L.~Leuzzi,
  Phys.\ Rev.\ A \textbf{91}, 053816 (2015).

\bibitem{Angelani2006}
  L.~Angelani, C.~Conti, G.~Ruocco, and F.~Zamponi,
  Phys.\ Rev.\ Lett.\ \textbf{96}, 065702 (2006).

\bibitem{Leuzzi2009}
  L.~Leuzzi, C.~Conti, V.~Folli, L.~Angelani, and G.~Ruocco,
  Phys.\ Rev.\ Lett.\ \textbf{101}, 107203 (2009).

\bibitem{Hackenbroich2003}
  G.~Hackenbroich, C.~Viviescas, and F.~Haake,
  Phys.\ Rev.\ A \textbf{68}, 063805 (2003).

\bibitem{Viviescas2003}
  C.~Viviescas and G.~Hackenbroich,
  Phys.\ Rev.\ A \textbf{67}, 013805 (2003).

\bibitem{Fallert2009}
  J.~Fallert et~al.,
  Nat.\ Photonics \textbf{3}, 279 (2009).

\bibitem{Conti2008}
  C.~Conti and A.~Fratalocchi,
  Nat.\ Phys.\ \textbf{4}, 794 (2008).

\bibitem{Zaitsev2010}
  O.~Zaitsev and L.~Deych,
  J.~Opt.\ \textbf{12}, 024001 (2010).

\bibitem{Marruzzo2018}
  A.~Marruzzo, P.~Tyagi, F.~Antenucci, A.~Pagnani, and L.~Leuzzi,
  SciPost Phys.\ \textbf{5}, 002 (2018).

\bibitem{Kirkpatrick1989}
  T.~R.~Kirkpatrick, D.~Thirumalai, and P.~G.~Wolynes,
  Phys.\ Rev.\ A \textbf{40}, 1045 (1989).

\bibitem{Marinari1994}
  E.~Marinari, G.~Parisi, and F.~Ritort,
  J.~Phys.\ A: Math.\ Gen.\ \textbf{27}, 7615 (1994).

\bibitem{Parisi1995}
  G.~Parisi and M.~Potters,
  J.~Phys.\ A: Math.\ Gen.\ \textbf{28}, 5267 (1995).

\bibitem{MPV1987}
  M.~M\'ezard, G.~Parisi, and M.~Virasoro,
  \textit{Spin Glass Theory and Beyond}
  (World Scientific, Singapore, 1987).

\bibitem{Viana1985}
  L.~Viana and A.~J.~Bray,
  J.~Phys.\ C \textbf{18}, 3037 (1985).

\end{thebibliography}

\end{document}
